% ===================================================================
% Front matter — Notation / List of Symbols
% Cambridge / Springer LNM convention.
% Monograph: Flos Aureus: A Theorem-Guided Architecture Search
%            through the Trinity Identity φ²+φ⁻²=3
% Author:    Кирилл Шамин (Cyril Shamin)
% Institute: University of Reading / gHashTag Open Doctorate
% ===================================================================
%
% Requirements: booktabs, longtable packages in preamble.
%
% TODO: This is a scaffold stub with ≥ 20 symbol rows.
%       Action the following before submission:
%
%   TODO-1: Add any chapter-local symbols not yet listed below.
%           Each symbol must be defined at its first use in the text
%           AND listed here — dual provenance is required (R4).
%   TODO-2: Verify all \ref{} cross-references compile without
%           warnings (run cargo run -p trios-phd -- compile).
%   TODO-3: If any constant is φ-derived (R6), annotate its row
%           with \textsf{[φ-derived]} in the Meaning column.
%   TODO-4: Confirm GF(16) row matches INV-3 in
%           assertions/igla_assertions.json byte-for-byte (R5).
%   TODO-5: Add missing operators (e.g., ∇, ⊗, ≅) as they appear
%           in chapters 10–16 (geometry / energy chapters).
% ===================================================================

\chapter*{Notation}
\addcontentsline{toc}{chapter}{Notation}
\thispagestyle{plain}

\noindent
The following table defines symbols used throughout the monograph.
Symbols local to a single section are defined at their point of first
use and are not repeated here. The column ``Ref.'' points to the
chapter or appendix where the symbol is introduced formally.

\bigskip

\begin{longtable}{@{}p{3.2cm}p{9cm}p{2cm}@{}}
\toprule
\textbf{Symbol} & \textbf{Meaning} & \textbf{Ref.} \\
\midrule
\endfirsthead
\toprule
\textbf{Symbol} & \textbf{Meaning} & \textbf{Ref.} \\
\midrule
\endhead
\bottomrule
\endfoot
%
% ── Golden-ratio family ──────────────────────────────────────────
\(\varphi\)
  & Golden ratio, \((1+\sqrt{5})/2 \approx 1.618\,033\,988\).
    \textsf{[φ-derived]}
  & Ch.~\ref{ch:golden-seed} \\
\(\psi\)
  & Algebraic conjugate of \(\varphi\):
    \(\psi = (1-\sqrt{5})/2 = -\varphi^{-1} \approx -0.618\).
  & Ch.~\ref{ch:golden-seed} \\
\(\varphi^{2}\)
  & \(\varphi+1 \approx 2.618\). \textsf{[φ-derived]}
  & Ch.~\ref{ch:golden-egg} \\
\(\varphi^{-2}\)
  & \(2-\varphi \approx 0.382\). \textsf{[φ-derived]}
  & Ch.~\ref{ch:golden-egg} \\
\(\varphi^{-3}\)
  & \(2\varphi-3 \approx 0.236\). Used as learning-rate seed
    \(\alpha_{\varphi} = \varphi^{-3}/2\). \textsf{[φ-derived]}
  & Ch.~\ref{ch:igla-race} \\
\(\alpha_{\varphi}\)
  & Learning-rate champion seed, \(\varphi^{-3}/2 \approx 0.118\).
    \textsf{[φ-derived]}
  & Ch.~\ref{ch:igla-race} \\
%
% ── Number sequences ─────────────────────────────────────────────
\(F_{n}\)
  & \(n\)-th Fibonacci number: \(F_{0}=0,\,F_{1}=1,\,F_{n}=F_{n-1}+F_{n-2}\).
  & Ch.~\ref{ch:fibonacci} \\
\(L_{n}\)
  & \(n\)-th Lucas number: \(L_{0}=2,\,L_{1}=1,\,L_{n}=L_{n-1}+L_{n-2}\).
  & Ch.~\ref{ch:lucas-ring} \\
\(\mathcal{L}\)
  & Lucas ring \(\mathbb{Z}[\varphi]\); the algebraic home of all
    φ-derived constants in this monograph.
  & Ch.~\ref{ch:lucas-ring} \\
%
% ── Algebra / geometry ───────────────────────────────────────────
\(\mathrm{GF}(16)\)
  & Galois field of order 16; precision floor of IGLA (INV-3).
  & Ch.~\ref{ch:experiments-gf16} \\
\(\langle \cdot,\cdot \rangle\)
  & Inner product in \(\mathbb{R}^{d}\) (Euclidean unless noted).
  & Ch.~\ref{ch:energy} \\
\(\|\cdot\|\)
  & Euclidean norm, \(\|x\|=\sqrt{\langle x,x\rangle}\).
  & Ch.~\ref{ch:energy} \\
\(\otimes\)
  & Kronecker / tensor product.
  & Ch.~\ref{ch:three-strands} \\
\(\cong\)
  & Ring or group isomorphism.
  & Ch.~\ref{ch:lucas-ring} \\
%
% ── Runtime invariants ───────────────────────────────────────────
\(\mathrm{INV}\text{-}k\)
  & The \(k\)-th IGLA runtime invariant, \(k=1,\ldots,5\);
    mechanized in \texttt{trinity-clara/proofs/igla/}.
    Full list in Appendix~F.
  & App.~F \\
%
% ── Empirical metrics ────────────────────────────────────────────
\(\mathrm{BPB}\)
  & Bits-per-byte; primary empirical loss metric for IGLA
    language models. Lower is better.
  & Ch.~\ref{ch:experiments-bpb} \\
\(\mathrm{ASHA}\)
  & Asynchronous Successive Halving Algorithm; hyperparameter
    scheduler used in all IGLA experiments.
  & Ch.~\ref{ch:experiments-asha} \\
\(\varepsilon\text{-greedy}\)
  & Exploration strategy with probability \(\varepsilon\) of random
    action; \(\varepsilon \in (0,1)\).
  & Ch.~\ref{ch:igla-race} \\
%
% ── Architecture components ──────────────────────────────────────
\(\mathrm{NCA}\)
  & Neural Cellular Automata; spatial-update architecture (INV-4).
  & Ch.~\ref{ch:nca} \\
\(\mathrm{JEPA}\)
  & Joint Embedding Predictive Architecture; an encoder family
    contrasted with IGLA in Chapter~\ref{ch:jepa}.
  & Ch.~\ref{ch:jepa} \\
\(\mathrm{VSA}\)
  & Vector Symbolic Architecture; holographic representation
    framework (INV-3, Ch.~\ref{ch:vsa}).
  & Ch.~\ref{ch:vsa} \\
%
% ── Proof conventions ────────────────────────────────────────────
\(\square\)
  & End of proof (open square = informal / sketch proof).
  & — \\
\(\blacksquare\)
  & End of proof (filled square = fully mechanized in Coq).
  & — \\
\(\partial\)
  & Boundary of a falsification region; see Appendix~B.
  & App.~B \\
%
% ── Notation TBD ─────────────────────────────────────────────────
% TODO-ADD: Insert additional symbols from chapters 10–16 here.
%
\end{longtable}
